% ============================================================
%  Diffusione Termica 1D — Fondamenti e Discretizzazione
%  Davide — Marzo 2026
% ============================================================
\documentclass[11pt, a4paper]{article}

% --- Pacchetti ---
\usepackage[T1]{fontenc}
\usepackage[utf8]{inputenc}
\usepackage[italian]{babel}
\usepackage{amsmath, amssymb}
\usepackage{physics}          % \dv, \pdv, ecc.
\usepackage{siunitx}
\usepackage{booktabs}
\usepackage{array}
\usepackage{geometry}
\usepackage{mathtools}
\usepackage{xcolor}
\usepackage{listings}
\usepackage{caption}
\usepackage{hyperref}
\usepackage{tikz}
\usetikzlibrary{arrows.meta, decorations.markings, patterns}

\geometry{margin=2.5cm}

\sisetup{
  output-decimal-marker = {,},
  per-mode = symbol
}

% Colori
\definecolor{myblue}{RGB}{30, 90, 180}
\definecolor{myred}{RGB}{180, 30, 30}
\definecolor{mygray}{RGB}{240, 240, 240}
\definecolor{codegreen}{RGB}{0,128,0}
\definecolor{codegray}{RGB}{128,128,128}

% Stile listing MATLAB
\lstdefinestyle{matlab}{
  language=Matlab,
  basicstyle=\ttfamily\small,
  keywordstyle=\color{myblue}\bfseries,
  commentstyle=\color{codegreen}\itshape,
  stringstyle=\color{myred},
  backgroundcolor=\color{mygray},
  frame=single,
  rulecolor=\color{gray!50},
  numbers=left,
  numberstyle=\tiny\color{codegray},
  numbersep=8pt,
  breaklines=true,
  captionpos=b
}

% Hyperref
\hypersetup{
  colorlinks=true,
  linkcolor=myblue,
  urlcolor=myblue
}

% ============================================================
\begin{document}

\begin{center}
  {\LARGE\bfseries\color{myblue} Diffusione Termica 1D}\\[6pt]
  {\large Fondamenti Fisici, Discretizzazione e Implementazione MATLAB}\\[8pt]
  {\normalsize Shannon \quad|\quad Marzo 2026}
\end{center}

\vspace{0.5em}
\hrule height 1.2pt
\vspace{1.5em}

\tableofcontents

\vspace{1.5em}
\hrule
\newpage

% ============================================================
\section{Geometria e Parametri del Problema}

Si considera una sbarra di alluminio con le seguenti caratteristiche geometriche e termofisiche:

\begin{table}[h!]
\centering
\caption{Parametri fisici del problema}
\begin{tabular}{llll}
\toprule
\textbf{Grandezza} & \textbf{Simbolo} & \textbf{Valore} & \textbf{Unità} \\
\midrule
Lunghezza                    & $L$       & \num{50}    & \si{\milli\meter} \\
Sezione trasversale          & $A$       & $\num{3} \times \num{30}$  & \si{\milli\meter\squared} \\
Diffusività termica (Al)     & $\alpha$  & \num{98e-6} & \si{\meter\squared\per\second} \\
Conducibilità termica (Al)   & $k$       & \num{237}   & \si{\watt\per\meter\per\kelvin} \\
Potenza imposta a $x=0$      & $q_\text{in}$ & \num{25} & \si{\watt} \\
Temperatura pozzo a $x=L$    & $T_\text{sink}$ & \num{80} & \si{\celsius} \\
Temperatura iniziale         & $T_0$     & \num{25}   & \si{\celsius} \\
\bottomrule
\end{tabular}
\end{table}

Il dominio spaziale è $x \in [0,\, L]$. Le condizioni al contorno sono:
\begin{itemize}
  \item \textbf{Bordo sinistro} ($x = 0$): flusso termico imposto \textbf{(Neumann)};
  \item \textbf{Bordo destro} ($x = L$): temperatura fissa \textbf{(Dirichlet)}.
\end{itemize}

\begin{figure}[h!]
\centering
\begin{tikzpicture}[scale=1.0]
  % Sbarra
  \fill[blue!10, draw=myblue, line width=1pt] (0,0) rectangle (8,1);
  % Freccia flusso entrante
  \draw[-{Latex[length=3mm]}, myred, line width=1.5pt] (-1.5, 0.5) -- (-0.05, 0.5);
  \node[myred, above] at (-0.75, 0.6) {$q_\text{in} = \SI{25}{\watt}$};
  % Etichetta bordo sinistro
  \node[myblue, below] at (0, -0.1) {$x=0$};
  \node[myblue, below] at (0, -0.55) {\small Neumann};
  % Pozzo termico a destra
  \fill[red!30, draw=myred, line width=1pt] (8,0) rectangle (9.2,1);
  \node[myred] at (8.6, 0.5) {$T_\text{sink}$};
  \node[myblue, below] at (8, -0.1) {$x=L$};
  \node[myblue, below] at (8, -0.55) {\small Dirichlet};
  % Asse x
  \draw[->, gray] (0,-0.9) -- (8,-0.9) node[right] {$x$};
  \node[gray, below] at (4,-1.15) {$L = \SI{50}{\milli\meter}$};
  % Freccia profilo temperatura
  \draw[->, dashed, myblue, line width=1pt] (0.3, 1.3) to[out=0, in=170] (7.7, 0.9);
  \node[myblue, above] at (4, 1.5) {\small Profilo $T(x,t) \to T_\text{ss}(x)$ lineare a regime};
\end{tikzpicture}
\caption{Schema del dominio: sbarra con flusso imposto a sinistra e pozzo termico a destra.}
\end{figure}

% ============================================================
\section{Equazione del Calore (PDE)}

\subsection{Forma Differenziale Generale}

L'equazione della diffusione del calore in un mezzo omogeneo e isotropo, in forma monodimensionale, è:

\begin{equation}
  \boxed{
    \frac{\partial T}{\partial t}(x,t) = \alpha \, \frac{\partial^2 T}{\partial x^2}(x,t)
  }
  \label{eq:heat}
\end{equation}

dove:
\begin{itemize}
  \item $T(x,t)$ è la distribuzione di temperatura \si{[\celsius]} in funzione della posizione $x$ e del tempo $t$;
  \item $\alpha = k/(\rho c_p)$ è la \emph{diffusività termica} \si{[\meter\squared\per\second]}, che misura la velocità con cui il calore si propaga nel materiale.
\end{itemize}

L'equazione~\eqref{eq:heat} discende dal bilancio energetico su un volume di controllo infinitesimo $A\,\dd x$, applicando la \textbf{legge di Fourier}:

\begin{equation}
  q = -k A \, \frac{\partial T}{\partial x}
  \label{eq:fourier}
\end{equation}

che afferma che il flusso di calore è proporzionale al gradiente di temperatura e orientato nella direzione del calore decrescente.

\subsection{Condizioni al Contorno}

\paragraph{BC sinistra — Neumann (flusso imposto).}
Il flusso netto uscente dalla sbarra nel verso $+x$ è nullo; la potenza $q_\text{in}$ entra dal bordo:

\begin{equation}
  -k A \, \frac{\partial T}{\partial x}\bigg|_{x=0} = q_\text{in}
  \quad \Longrightarrow \quad
  \frac{\partial T}{\partial x}\bigg|_{x=0} = -\frac{q_\text{in}}{k A}
  \label{eq:BC_neumann}
\end{equation}

\paragraph{BC destra — Dirichlet (temperatura imposta).}

\begin{equation}
  T(L, t) = T_\text{sink} = \SI{80}{\celsius} \qquad \forall\, t > 0
  \label{eq:BC_dirichlet}
\end{equation}

\subsection{Condizione Iniziale}

La sbarra si trova inizialmente a temperatura uniforme:

\begin{equation}
  T(x, 0) = T_0 = \SI{25}{\celsius} \qquad \forall\, x \in [0, L]
  \label{eq:IC}
\end{equation}

% ============================================================
\section{Soluzione a Regime (Stazionaria)}

A regime ($\partial T / \partial t = 0$) l'equazione~\eqref{eq:heat} si riduce a:

\begin{equation}
  \frac{\dd^2 T}{\dd x^2} = 0
  \label{eq:laplace1D}
\end{equation}

La soluzione generale è un profilo \textbf{lineare}:

\begin{equation}
  T_\text{ss}(x) = a\, x + b
  \label{eq:lin}
\end{equation}

Applicando le condizioni~\eqref{eq:BC_neumann} e~\eqref{eq:BC_dirichlet}:

\begin{align}
  a &= \frac{\dd T_\text{ss}}{\dd x} = -\frac{q_\text{in}}{k A} \label{eq:a}\\[6pt]
  b &= T_\text{sink} - a\, L \label{eq:b}
\end{align}

Sostituendo i valori numerici:

\begin{equation}
  a = -\frac{\SI{25}{\watt}}{\SI{237}{\watt\per\meter\per\kelvin} \times \SI{9e-5}{\meter\squared}}
    = \SI{-1170.96}{\kelvin\per\meter}
\end{equation}

\begin{equation}
  T_\text{ss}(0) = T_\text{sink} - a\,L = 80 - (-1170.96) \times 0.05 \approx \SI{138.5}{\celsius}
\end{equation}

La temperatura al bordo caldo è quindi circa \SI{138.5}{\celsius}: la caduta di \SI{58.5}{\celsius} lungo i \SI{50}{\milli\meter} è dovuta interamente alla resistenza termica della sbarra.

% ============================================================
\section{Discretizzazione: Metodo delle Differenze Finite (FTCS)}

\subsection{Griglia Spazio-Temporale}

Il dominio continuo $[0,L] \times [0, T_\text{final}]$ viene discretizzato in:
\begin{itemize}
  \item $N_x$ nodi spaziali con passo $\Delta x = L / (N_x - 1)$;
  \item passi temporali di ampiezza $\Delta t$ (determinati dalla condizione di stabilità).
\end{itemize}

Si indica con $u_i^n \approx T(i\,\Delta x,\, n\,\Delta t)$ il valore discreto della temperatura.

\subsection{Schema FTCS per i Nodi Interni}

Si applica lo schema \textbf{Forward-Time, Centered-Space (FTCS)}: derivata temporale con differenza in avanti, derivata spaziale seconda con differenza centrata:

\begin{equation}
  \frac{u_i^{n+1} - u_i^n}{\Delta t} = \alpha \, \frac{u_{i+1}^n - 2u_i^n + u_{i-1}^n}{\Delta x^2}
  \label{eq:ftcs}
\end{equation}

Raccogliendo, si ottiene la formula esplicita di aggiornamento:

\begin{equation}
  \boxed{
    u_i^{n+1} = u_i^n + \lambda \left( u_{i+1}^n - 2u_i^n + u_{i-1}^n \right)
  }
  \quad i = 2, \ldots, N_x - 1
  \label{eq:ftcs_explicit}
\end{equation}

dove $\lambda$ è il \textbf{numero di Fourier discreto} (o coefficiente CFL):

\begin{equation}
  \lambda = \frac{\alpha\, \Delta t}{\Delta x^2}
  \label{eq:lambda}
\end{equation}

\subsection{Condizione di Stabilità di Von Neumann}

Lo schema FTCS è \emph{condizionatamente stabile}. L'analisi di Von Neumann richiede:

\begin{equation}
  \lambda \leq \frac{1}{2}
  \label{eq:stability}
\end{equation}

Nella simulazione si sceglie $\lambda = 0{,}45$ (safety factor), da cui si ricava:

\begin{equation}
  \Delta t = \frac{0{,}45 \cdot \Delta x^2}{\alpha}
  \label{eq:dt}
\end{equation}

Se $\lambda > 0{,}5$ la soluzione numerica \emph{diverge}, indipendentemente dalla qualità della griglia spaziale.

% ============================================================
\section{Implementazione delle Condizioni al Contorno}

\subsection{BC Destra — Dirichlet (nodo $N_x$)}

La condizione più semplice: il nodo all'estremo destro è vincolato al valore fisso ad ogni passo:

\begin{equation}
  u_{N_x}^{n+1} = T_\text{sink} \qquad \forall\, n
  \label{eq:dirichlet_impl}
\end{equation}

\subsection{BC Sinistra — Neumann con Ghost Node (nodo $i=1$)}

La formula FTCS~\eqref{eq:ftcs_explicit} richiede il nodo $u_0^n$ (a sinistra del bordo), che non esiste fisicamente. Si introduce un \textbf{ghost node} virtuale $u_0^n$, determinato imponendo la derivata al bordo con una differenza centrata:

\begin{equation}
  \frac{\partial T}{\partial x}\bigg|_{x=0} \approx \frac{u_2^n - u_0^n}{2\,\Delta x} = -\frac{q_\text{in}}{k A}
\end{equation}

Da cui:

\begin{equation}
  u_0^n = u_2^n + \frac{2\,\Delta x\, q_\text{in}}{k A}
  \label{eq:ghost}
\end{equation}

Sostituendo~\eqref{eq:ghost} nell'equazione di aggiornamento per $i=1$:

\begin{align}
  u_1^{n+1} &= u_1^n + \lambda \left( u_2^n - 2u_1^n + u_0^n \right) \notag\\
             &= u_1^n + \lambda \left( u_2^n - 2u_1^n + u_2^n + \frac{2\,\Delta x\, q_\text{in}}{k A} \right) \notag\\
  \Aboxed{u_1^{n+1} &= u_1^n + \lambda \left( 2u_2^n - 2u_1^n + \frac{2\,\Delta x\, q_\text{in}}{k A} \right)}
  \label{eq:neumann_impl}
\end{align}

Questo approccio garantisce \textbf{ordine 2} nell'approssimazione spaziale della BC, a differenza di una semplice differenza in avanti che sarebbe solo del primo ordine.

% ============================================================
\section{Riepilogo del Loop Temporale}

Ad ogni passo $n \to n+1$ la simulazione esegue, nell'ordine:

\begin{enumerate}
  \item \textbf{Aggiornamento nodi interni} ($i = 2, \ldots, N_x-1$) con la formula~\eqref{eq:ftcs_explicit};
  \item \textbf{BC sinistra} — nodo $i=1$ con la formula~\eqref{eq:neumann_impl} (ghost node);
  \item \textbf{BC destra} — nodo $i=N_x$ fissato a $T_\text{sink}$ per la Dirichlet~\eqref{eq:dirichlet_impl}.
\end{enumerate}

In MATLAB, i passi 1--2 si implementano in forma vettorizzata (senza loop espliciti), garantendo l'efficienza tipica di MATLAB:

\begin{lstlisting}[style=matlab, caption={Kernel del loop temporale in MATLAB}]
% Nodi interni (FTCS)
u_new(2:end-1) = u(2:end-1) + lambda * ...
    (u(3:end) - 2*u(2:end-1) + u(1:end-2));

% BC sinistra: ghost node Neumann
u_ghost  = u(2) + 2 * dx * q_in / (k * A);
u_new(1) = u(1) + lambda * (u(2) - 2*u(1) + u_ghost);

% BC destra: Dirichlet
u_new(end) = T_sink;
\end{lstlisting}

% ============================================================
\section{Verifica: Errore a Regime}

Al termine della simulazione, si confronta la soluzione numerica $u_i^{N_t}$ con il profilo analitico $T_\text{ss}(x_i)$:

\begin{equation}
  \varepsilon_\text{max} = \max_{i} \left| u_i^{N_t} - T_\text{ss}(x_i) \right|
\end{equation}

Con $N_x = 200$ e $\lambda = 0{,}45$, ci si aspetta $\varepsilon_\text{max} \ll \SI{0.1}{\celsius}$, confermando la correttezza della discretizzazione e la convergenza al regime fisico previsto.

% ============================================================
\appendix
\section{Nomenclatura}

\begin{tabular}{>{\bfseries}ll}
$T(x,t)$         & Temperatura \si{[\celsius]}\\
$\alpha$         & Diffusività termica \si{[\meter\squared\per\second]}\\
$k$              & Conducibilità termica \si{[\watt\per\meter\per\kelvin]}\\
$A$              & Area della sezione trasversale \si{[\meter\squared]}\\
$L$              & Lunghezza della sbarra \si{[\meter]}\\
$q_\text{in}$    & Potenza termica imposta \si{[\watt]}\\
$T_\text{sink}$  & Temperatura del pozzo termico \si{[\celsius]}\\
$\lambda$        & Numero di Fourier discreto (coefficiente CFL) \si{[-]}\\
$\Delta x$       & Passo spaziale \si{[\meter]}\\
$\Delta t$       & Passo temporale \si{[\second]}\\
$u_i^n$          & Temperatura discreta al nodo $i$, istante $n$\\
$u_0^n$          & Ghost node virtuale per la BC di Neumann\\
\end{tabular}

\end{document}
